% SPDX-License-Identifier: CC-BY-SA-4.0
% Author: Matthieu Perrin
% Part: 
% Section: 
% Sub-section: 
% Frame: 

\begingroup

\begin{frame}[fragile]{Progrès bloquant versus progrès non-bloquant}

  \begin{description}[Progrès non-bloquant :]
  \item[Progrès non-bloquant :] Les threads corrects progressent \alert{malgré les pannes}.
    \begin{itemize}
    \item propriété violée par l'utilisation de verrous.
    \end{itemize}
  \item[Progrès bloquant :] Les threads progressent \alert{s'il n'y a pas de panne}.
    \begin{itemize}
    \item propriété violée lors d'un interblocage circulaire.
    \end{itemize}
  \end{description}
  
  \begin{block}{Définition plus précise du progrès bloquant}
    Un objet $O$ vérifie le progrès bloquant si,
    pour toute exécution dans laquelle aucun processus ne tombe en panne,
    toutes les opérations se terminent.
  \end{block}

  \begin{exampleblock}{Exemple d'algorithme bloquant}
    \begin{lstlisting}[gobble=6]
      int getAndIncrement() {
        while (lock.getAndSet(true));
        value.set(value.get()+1);
        lock.set(false);
      }
    \end{lstlisting}
  \end{exampleblock}

\end{frame}

\endgroup
\endinput
